% Custom beamer theme for the CGCNN presentation.
% ================================================
%
% Kept separate from the content file so the visual system - colours, fonts,
% the card/eyebrow building blocks - can be read and adjusted in one place
% without wading through forty frames of prose.
%
% DESIGN
% ------
% Georgia for headings (a serif with the same weight/character as the
% Cambria used in the project's other teaching decks), Avenir Next for body
% text (same role as their Calibri). Both are macOS system fonts, loaded via
% fontspec, so this MUST be built with xelatex or lualatex - not pdflatex.
%
% Cards are TikZ nodes with `text width` set and NO fixed height. That one
% choice eliminates an entire class of bug: a fixed-height text box that
% silently overflows its frame is exactly what an earlier (PowerPoint)
% version of this deck kept getting wrong. A TikZ node sized only by its
% content cannot overflow itself - it grows to fit, by construction.

\usepackage{fontspec}
\setmainfont{Avenir Next}
\setsansfont{Avenir Next}
\newfontfamily\headingfont{Georgia}
\newfontfamily\monofont{Menlo}

\usepackage{tikz}
\usetikzlibrary{arrows.meta,positioning,shapes.geometric,calc,backgrounds,fit}
\usepackage{amsmath,amssymb}
\usepackage{graphicx}
\usepackage{booktabs}
% Both results/ subdirectories are searched in order - every existing
% \includegraphics{parity_K_VRH_ens} etc. call (bare filename, no path
% prefix) keeps resolving with no per-call edits after the cgcnn/alignn split.
\graphicspath{{figures/}{../results/cgcnn/}{../results/alignn/}}

% --- Palette ----------------------------------------------------------------
\definecolor{Navy}{HTML}{0F1B33}
\definecolor{SlideBlue}{HTML}{1450AA}
\definecolor{LightBlue}{HTML}{6FA8FF}
\definecolor{SlideGreen}{HTML}{1E8E5A}
\definecolor{SlideAmber}{HTML}{D97B12}
\definecolor{SlideRed}{HTML}{B02A2A}
\definecolor{SlideInk}{HTML}{1F2937}
\definecolor{Muted}{HTML}{5B6B82}
\definecolor{Subtle}{HTML}{AAB6C8}
\definecolor{CardBg}{HTML}{F4F6FA}
\definecolor{CardRule}{HTML}{D8E0EC}
\definecolor{GreenBg}{HTML}{EFF5EF}
\definecolor{GreenRule}{HTML}{C8E0C8}
\definecolor{AmberBg}{HTML}{FDF4E8}
\definecolor{AmberRule}{HTML}{F0D8B8}
\definecolor{BlueBg}{HTML}{E8EFFA}
\definecolor{RedBg}{HTML}{FDF0F0}
\definecolor{RedRule}{HTML}{F0C8C8}
\definecolor{NaColor}{HTML}{7A3FA0}   % matches figures/nacl_*.png
\definecolor{ClColor}{HTML}{1E8E5A}

% --- Bare beamer plumbing: no default chrome, we draw our own -------------
\setbeamertemplate{navigation symbols}{}
\setbeamertemplate{footline}{%
  \hfill\color{Subtle}\scriptsize\insertframenumber\,/\,\inserttotalframenumber\hspace{4mm}\vspace{2mm}
}
\setbeamercolor{background canvas}{bg=white}
\setbeamersize{text margin left=0.85cm, text margin right=0.85cm}
\setbeamertemplate{itemize item}{\textcolor{SlideBlue}{\raisebox{0.15em}{\rule{5pt}{5pt}}}}
\setbeamertemplate{itemize subitem}{\textcolor{Muted}{--}}
\setlength{\parskip}{0pt}

% frametitle is unused throughout - every content frame builds its own header
% with \SlideHeader, so its template is blanked out rather than left to draw
% something nobody positioned.
\setbeamertemplate{frametitle}{}
\setbeamertemplate{headline}{}

% --- The eyebrow + title header every content slide starts with -----------
\newcommand{\Chapter}[0]{}
\newcommand{\SlideHeader}[2]{%
  \par\noindent{\color{SlideBlue}\bfseries\sffamily\fontsize{10}{12}\selectfont #1}\\[2pt]
  \noindent{\color{SlideInk}\bfseries\headingfont\fontsize{19}{22}\selectfont #2}\\[5pt]%
}

% --- A soft card panel; sizes to its content, never the other way round ----
% #1 width  #2 fill  #3 rule (use `none` for no border)  #4 content
\newcommand{\Card}[4]{%
  \begin{tikzpicture}
    \node[rounded corners=4pt, fill=#2,
          draw=#3, line width=0.7pt,
          text width=#1, inner sep=6pt, align=left,
          minimum height=0pt] {#4};
  \end{tikzpicture}%
}
% A card with a coloured accent bar down the left edge - used for the three-
% way "why" explainer boxes.
\newcommand{\AccentCard}[3]{%
  \begin{tikzpicture}
    \node[rounded corners=4pt, fill=CardBg, draw=CardRule, line width=0.7pt,
          text width=#1, inner sep=6pt, align=left,
          minimum height=0pt] (box) {#3};
    \fill[#2, rounded corners=2pt] (box.north west) rectangle
      ($(box.south west)+(0.09,0)$);
  \end{tikzpicture}%
}

% --- The "so what" strip anchored to the bottom of a frame -----------------
\newcommand{\Takeaway}[1]{%
  \vfill
  \begin{tikzpicture}
    \node[rounded corners=4pt, fill=CardBg, draw=none,
          text width=11.3cm, inner sep=6pt, align=left,
          minimum height=0pt] (box) {\sffamily\bfseries\footnotesize\color{SlideInk} #1};
    \fill[SlideBlue, rounded corners=2pt] (box.north west) rectangle
      ($(box.south west)+(0.08,0)$);
  \end{tikzpicture}%
}

% --- Section-divider frame: full navy, one big line of text ---------------
\newenvironment{dividerframe}{%
  \begingroup
  \setbeamercolor{background canvas}{bg=Navy}
  \begin{frame}[plain]
  \vspace{2.6cm}
}{%
  \end{frame}
  \endgroup
}

% --- Title-slide flow pill (the row of stage names on slide 1) ------------
\newcommand{\FlowPill}[1]{%
  \tikz\node[rounded corners=4pt, fill=Navy!60, draw=LightBlue!40,
             line width=0.5pt, text width=2.1cm, align=center,
             inner sep=7pt, font=\sffamily\bfseries\small\color{white}]
             {#1};%
}

\input{metrics.tex}
